%!TEX TS-program = pdflatex
%!TEX TS-options = -shell-escape
% Die Lösungsschalter bleiben wie bei den anderen Blättern in Zeile 4 und 5.
\newcommand{\printpraesenzlsg}{false}
\newcommand{\printloesungen}{true}
\newcommand{\printbewertungen}{false}
% Für die Lösungsausgabe in Zeile 5 false durch true ersetzen.
\newcommand{\blattnummer}{2}
\documentclass[%
  paper=a4,
  fontsize=10pt,
  ngerman
  ]{scrartcl}

% Basics für Codierung und Sprache
% ===========================================================
  \usepackage{scrtime}
  \usepackage{etex}
  \usepackage{shellesc}
  \usepackage[final]{graphicx}
  \usepackage[utf8]{inputenc}
  \usepackage{babel}
  \usepackage[german=quotes]{csquotes}
  \usepackage{datetime2}
  \usepackage{ifthen}
  \usepackage{dingbat}
% ===========================================================

%% Fonts und Typographie
%% ===========================================================
  \usepackage{charter}
  \usepackage{nimbusmononarrow}
  \usepackage[babel=true,final,tracking=smallcaps]{microtype}
  \DisableLigatures{encoding = T1, family = tt* }
  \usepackage{ellipsis}
  \usepackage{listings}
  
% ===========================================================

% Farben
% ===========================================================
  \usepackage[usenames,x11names,table]{xcolor}
  \definecolor{fbblau}{HTML}{3078AB}
  \definecolor{blackberry}{rgb}{0.53, 0.0, 0.25}
% ===========================================================

% Mathe-Pakete und -Einstellungen
% ===========================================================
  \usepackage{mathtools}
  \usepackage{amssymb}
  \usepackage[bigdelims]{newtxmath}
  \allowdisplaybreaks
  \usepackage{bm}
  \usepackage{wasysym}
% ===========================================================

% TikZ
% ===========================================================
  \usepackage{tikz}
  \usepackage{tikz-cd}
  \usepackage{pgf}
  \usetikzlibrary{matrix}
  \usetikzlibrary{positioning}
  \usetikzlibrary{arrows}
  \usetikzlibrary{shapes}
  \tikzset{>=Latex}
  \usepackage[mode=image]{standalone}
% ===========================================================

% Seitenlayout, Kopf-/Fußzeile
% ===========================================================
  \usepackage{scrlayer-scrpage}
  \usepackage[top=3cm, bottom=2.5cm, left=2.5cm, right=2cm]{geometry}  
  \pagestyle{scrheadings}
  \clearscrheadfoot 
  \setheadsepline{0.4pt}
  \setfootsepline{0.4pt}
  \setkomafont{pagehead}{\normalfont\footnotesize}
  \setkomafont{pagefoot}{\normalfont\footnotesize}
  \lohead[]{Übung zur Vorlesung \textit{Embedded Systems Engineering} -- Wintersemester 2026}
  \rohead[]{%
    \ifthenelse{\boolean{loesungen}}{\textbf{Lösungen zu Blatt \blattnummer}}{Blatt \blattnummer}}
  %\rofoot[Version: \DTMnow]{%
  %  \ifthenelse{\boolean{loesungen}}{Version: \DTMnow}{}}
  %\lofoot[\jobname.tex]{%
  %  \ifthenelse{\boolean{loesungen}}{\jobname.tex}{}}
  \cfoot[\thepage]{\thepage}
  \raggedbottom
  \usepackage{setspace}          
  \linespread{1.2}
  \setlength{\parindent}{0pt}
  \setlength{\parskip}{0.5\baselineskip}
  \usepackage[all]{nowidow}
  \usepackage{lscape}
% ===========================================================

% Hyperref
% ===========================================================
  \usepackage[%
    hidelinks,
    pdfpagelabels,
    bookmarksopen=true,
    bookmarksnumbered=true,
    linkcolor=black,
    urlcolor=fbblau,
    plainpages=false,
    pagebackref,
    citecolor=black,
    hypertexnames=true,
    pdfborderstyle={/S/U},
    linkbordercolor=fbblau,
    colorlinks=false,
    backref=false]{hyperref}
  \hypersetup{final}
  \makeatletter
  \g@addto@macro{\UrlBreaks}{\UrlOrds}
  \makeatother
% ===========================================================

% Listen und Tabellen
% ===========================================================
  \usepackage{multicol}
  \usepackage{multirow}
  \usepackage[shortlabels]{enumitem}
  \setlist[enumerate]{font=\bfseries,itemsep=0.25\baselineskip}
  \setlist[itemize]{label=$\triangleright$,itemsep=-0.25\baselineskip}
  \usepackage{tabularx}
  \usepackage{array}
  \usepackage{listings}
  \lstset{breaklines=true,
    basicstyle=\ttfamily,
    frame=tblr,
    language=c,
    numbers=left}
  \usepackage{include/mips}

\definecolor{lightgreen}{rgb}{0,0.7,0}
\definecolor{gray}{rgb}{0.5,0.5,0.5}
\definecolor{mauve}{rgb}{0.58,0,0.82}

\lstdefinestyle{mips}{ %
  language=[mips]Assembler,       % the language of the code
  keywordstyle=\color{blue},          % keyword style
  keywordstyle=[2]\color{magenta},     % keyword style
  keywordstyle=[3]\color{red},          % keyword style
  commentstyle=\color{lightgreen},       % comment style
  stringstyle=\color{mauve},         % string literal style
  escapeinside={\%*}{*)},            % if you want to add a comment within your code
  morekeywords={*,...}               % if you want to add more keywords to the set
}

\newcommand{\mipsinline}[1]{\lstinline[style=mips,columns=fixed]{#1}}
\newcommand{\mipsregister}[1]{\mipsinline{\$#1}}
  \newcommand{\qed}{\hfill \ensuremath{\Box}}
  \newcommand{\cmd}[1]{\texttt{#1}}
% ===========================================================

% Sonstiges
% ===========================================================
  \usepackage[textsize=small,color=Red1!80!OrangeRed1!80]{todonotes}
  \usepackage{lipsum}
% ===========================================================

% Shortcuts and Math Commands
% ===========================================================
  \newcommand{\BB}{\mathbb{B}}      % Bool
  \newcommand{\CC}{\mathbb{C}}      % komplexe Zahlen
  \newcommand{\NN}{\mathbb{N}}      % nat. Zahlen
  \newcommand{\QQ}{\mathbb{Q}}      % rat. Zahlen
  \newcommand{\RR}{\mathbb{R}}      % reelle Zahlen
  \newcommand{\ZZ}{\mathbb{Z}}      % ganze Zahlen
  \newcommand{\oh}{\mathcal{O}}      % Landau-O
  \newcommand{\ol}[1]{\overline{#1}}    % überstreichen
  \newcommand{\wt}[1]{\widetilde{#1}}    % überschlängeln
  \newcommand{\wh}[1]{\widehat{#1}}    % überdachen
  \newcommand{\entspricht}{\mathrel{\widehat{=}}}
  
  \DeclarePairedDelimiter{\absolut}{\lvert}{\rvert}    % Betrag
  \DeclarePairedDelimiter{\ceiling}{\lceil}{\rceil}    % aufrunden
  \DeclarePairedDelimiter{\Floor}{\lfloor}{\rfloor}    % aufrunden
  \DeclarePairedDelimiter{\Norm}{\lVert}{\rVert}      % Norm
  \DeclarePairedDelimiter{\sprod}{\langle}{\rangle}    % spitze Klammern
  \DeclarePairedDelimiter{\enbrace}{(}{)}          % runde Klammern
  \DeclarePairedDelimiter{\benbrace}{\lbrack}{\rbrack}  % eckige Klammern
  \DeclarePairedDelimiter{\penbrace}{\{}{\}}        % geschweifte Klammern
  \newcommand{\Underbrace}[2]{{\underbrace{#1}_{#2}}}   % bessere Unterklammerungen
  
  % Kurzschreibweisen für Faule und Code-Vervollständigung
  \newcommand{\abs}[1]{\absolut*{#1}}
  \newcommand{\ceil}[1]{\ceiling*{#1}}
  \newcommand{\flo}[1]{\Floor*{#1}}
  \newcommand{\no}[1]{\Norm*{#1}}
  \newcommand{\sk}[1]{\sprod*{#1}}
  \newcommand{\enb}[1]{\enbrace*{#1}}
  \newcommand{\penb}[1]{\penbrace*{#1}}
  \newcommand{\benb}[1]{\benbrace*{#1}}
  \newcommand{\stack}[2]{\makebox[1cm][c]{$\stackrel{#1}{#2}$}}

% Übungszettel-Krams
% ===========================================================
  \usepackage[tikz]{mdframed}
  \newmdenv[%
    backgroundcolor = gray!20,
    linewidth=0pt,
    skipabove = 0.5cm,
    skipbelow = 0cm
  ]{notes}
  
  % Lösungsausgabe
  \newboolean{praesenzlsgn}
  \setboolean{praesenzlsgn}{\printpraesenzlsg}
  \newboolean{loesungen}
  \setboolean{loesungen}{\printloesungen}
  % Bewertungsschema-Ausgabe
  \newboolean{bewertungen}
  \setboolean{bewertungen}{\printbewertungen}

  \newcommand{\iforiginal}[1]{\ifthenelse{\boolean{praesenzlsgn}}{}{#1}}
  \newcommand{\ifpraesenzlsg}[1]{\ifthenelse{\boolean{praesenzlsgn}}{#1}{}}
  \newcommand{\ifloesung}[1]{\ifthenelse{\boolean{loesungen}}{#1}{}}
  \newcommand{\ifbewertung}[1]{\ifthenelse{\boolean{bewertungen}}{#1}{}}

  \newcommand{\umbruchoriginal}{\iforiginal{\newpage}}        % Seitenumbruch nur im Original
  \newcommand{\umbruchpraesenzlsg}{\ifpraesenzlsg{\newpage}}  % Seitenumbruch nur in Lösung d. Präsenzaufgabe
  \newcommand{\umbruchlsg}{\ifloesung{\newpage}}              % Seitenumbruch nur in Lösung
  \newcommand{\umbruchbewertung}{\ifbewertung{\newpage}}      % Seitenumbruch nur in Lösung mit Bewertungsschema
  \newcommand{\umbruchnurpraesenz}{\ifthenelse{\boolean{loesungen}}{}{\umbruchpraesenzlsg}}
  \newcommand{\umbruchnurlsg}{\ifthenelse{\boolean{bewertungen}}{}{\umbruchlsg}}

  \usepackage{comment}
  \newenvironment{praesenzlsg}%
    {\ifthenelse{\boolean{praesenzlsgn}}{\textbf{--- Präsenzaufgabe Anfang ---} \mbox{} \newline}{}}%
	{\ifthenelse{\boolean{praesenzlsgn}}{\mbox{} \newline \textbf{--- Präsenzaufgabe Ende ---}}{}}
  \newenvironment{loesung}%
    {\ifthenelse{\boolean{loesungen}}{\textbf{--- Lösung Anfang ---} \mbox{} \newline}{}}%
    {\ifthenelse{\boolean{loesungen}}{\mbox{} \newline \textbf{--- Lösung Ende ---} \ifthenelse{\boolean{bewertungen}}{}{\umbruchlsg}}{}}
  \newenvironment{bewertung}%
    {\ifthenelse{\boolean{bewertungen}}{\textbf{--- Bewertungsschema Anfang ---} \mbox{} \newline}{}}%
    {\ifthenelse{\boolean{bewertungen}}{\mbox{} \newline \textbf{--- Bewertungsschema Ende ---} \umbruchlsg}{}}

  \ifthenelse{\boolean{praesenzlsgn}}{}{\excludecomment{praesenzlsg}}
  \ifthenelse{\boolean{loesungen}}{}{\excludecomment{loesung}}
  \ifthenelse{\boolean{bewertungen}}{}{\excludecomment{bewertung}}

  % Aufgabennummerierung
  \newcounter{aufgabennummer} 
  \setcounter{aufgabennummer}{1}
  \newcommand{\aufgabe}[1]{\textbf{Aufgabe \blattnummer.\theaufgabennummer} \stepcounter{aufgabennummer} \hfill (#1 Punkte)}
  %\newcommand{\aufgabetitel}[2]{\iforiginal{\vspace{0.5cm}} \textbf{Aufgabe \blattnummer.\theaufgabennummer \ (#2)} \stepcounter{aufgabennummer} \hfill (#1 Punkte)}
  \newcommand{\aufgabetitel}[2]{\iforiginal{\vspace{0.5cm}} \textbf{Aufgabe \theaufgabennummer \ (#2)} \stepcounter{aufgabennummer} \hfill (#1 Punkte)}
  \newcommand{\praesenztitel}[1]{\iforiginal{\vspace{0.5cm}} \textbf{Aufgabe \blattnummer.\theaufgabennummer \ (#1)} \stepcounter{aufgabennummer} \hfill \smallpencil}

\newcommand{\danger}{\raisebox{0.5mm}{\fontencoding{U}\fontfamily{futs}\selectfont\char 66\relax}}
\newcommand{\ddanger}[1]{{\danger} \textbf{#1} {\danger}}

% Blatt-Thema (optional)
\newcommand{\blatthema}{} % default: leer
\newcommand{\setblatthema}[1]{\renewcommand{\blatthema}{#1}}

% Rechnerstrukturen-Krams
% ===========================================================
  \newcommand{\Band}{\ensuremath{\mathbin{\texttt{AND}}}}
  \newcommand{\Bor}{\ensuremath{\mathbin{\texttt{OR}}}}
  \newcommand{\Bnot}{\ensuremath{\mathbin{\texttt{NOT}}}}
  \newcommand{\Bnand}{\ensuremath{\mathbin{\texttt{NAND}}}}
  \newcommand{\Bnor}{\ensuremath{\mathbin{\texttt{NOR}}}}
  \newcommand{\Bxor}{\ensuremath{\mathbin{\texttt{XOR}}}}
  \newcommand{\Bxnor}{\ensuremath{\mathbin{\texttt{XNOR}}}}
  \newcommand{\Bone}{\ensuremath{\mathbin{\texttt{ONE}}}}
  \newcommand{\Bzero}{\ensuremath{\mathbin{\texttt{ZERO}}}}
  \newcommand{\Binary}[1]{\ensuremath{\texttt{#1}}\xspace}
  \newcommand{\ieee}[3]{\fbox{\ensuremath{\texttt{#1}}}\fbox{\ensuremath{\texttt{#2}}}\fbox{\ensuremath{\texttt{#3}}}}
  \DeclareMathOperator{\DNF}{DNF}
  \DeclareMathOperator{\KNF}{KNF}
  \DeclareMathOperator{\xor}{\oplus}
  \newcommand{\dc}{\star} %don't-care-Wert


% Create Code Listings
% ===========================================================
  \lstset{
  basicstyle=\ttfamily,
  numbers=left,
  numberstyle=\small,
  stepnumber=1,
  numbersep=10pt,
  frame=single,
  columns=fullflexible,
  keepspaces=true,
  showstringspaces=false,
  xleftmargin=1.5em
}
 
\usetikzlibrary{arrows.meta}
\tikzset{
  pnplace/.style={circle,draw,minimum size=8mm,inner sep=0pt},
  pntrans/.style={rectangle,draw,fill=gray!15,minimum width=5mm,minimum height=9mm,inner sep=0pt},
  pnlabel/.style={font=\small,align=center},
  pnarc/.style={->,>=Latex,semithick}
}
\newcommand{\pnmark}[1]{\fill (#1.center) circle (2pt);}

\begin{document}
\setblatthema{Petrinetze: Einfache Modellierungen}
\iforiginal{\thispagestyle{scrplain}
\vspace*{-3cm}
\begin{minipage}[t][1.1cm][c]{4.5cm}
  \includegraphics[width=4cm]{include/hwr-logo.png}
\end{minipage}
\hfill
\begin{minipage}[t][1.5cm][c]{8cm}
  \begin{center}
  \begin{footnotesize}
    \textsf{Prof. Dr.-Ing. Sebastian Schlesinger} \\[-0.1cm]
    \textsf{Fachbereich 2 - Duales Studium Wirtschaft \& Technik}
    %\url{https://www.hwr-berlin.de}
  \end{footnotesize}
  \end{center}
\end{minipage}
\hfill
%\begin{minipage}[t][1.6cm][c]{3.2cm}
%  \includegraphics[width=2.8cm]{include/ESlogo2}
%\end{minipage}

\vspace*{-0.3cm}
\begin{center} 
  \hrulefill \\[0.1cm]
  {\large Übung} \\[0.15cm]
    {\huge \bfseries Embedded Systems Engineering} \\[0.10cm]
  \ifthenelse{\equal{\blatthema}{}}{}{%
    {\Large \bfseries \blatthema} \\[0.10cm]
  }%
  {HWR Berlin, Wintersemester 2026} \\[-0.4cm]
  \begin{tabular}{lcr}
    \hspace{0.3\textwidth}   & \hspace{0.3\textwidth} & \hspace{0.3\textwidth} \\
    Prof. Dr.-Ing. Sebastian Schlesinger   &                        & "Ubung    \\ 
    Blatt \blattnummer     % \abgabetermin          \\ 
  \end{tabular} \\[0.1cm]
  \hrulefill
\end{center}}

\begin{notes}
\textbf{Grundlage:} Vorlesung ,,Embedded Systems -- Modellierung``, Folien 63--65.\par
Modelliert die Abläufe als Petrinetze und erläutert ihr Verhalten direkt an der Zeichnung.
Stellen werden als Kreise, Transitionen als Rechtecke und Marken als Punkte gezeichnet.
Verbindet durch Pfeile immer eine Stelle mit einer Transition oder umgekehrt.
Alle Kanten haben Gewicht eins.
Eine Transition ist aktiviert, wenn jede Eingangsstelle mindestens eine Marke enthält.
Beim Schalten wird je eine Marke aus jeder Eingangsstelle entfernt und in jeder
Ausgangsstelle erzeugt. Es gibt keine zusätzlichen Kapazitätsgrenzen oder Zeitbedingungen.
Beschriftet Stellen und Transitionen mit ihrer Bedeutung und zeichnet die Marken
im Anfangszustand ein. Eine formale Beschreibung durch Mengen oder Formeln ist nicht erforderlich.
\end{notes}

\aufgabetitel{$5$}{Bearbeitungsstation}\\
Eine Station bearbeitet jeweils ein Werkstück. Anfangs ist sie frei.
Durch das Einlegen eines Werkstücks beginnt die Bearbeitung. Nach deren Abschluss
liegt das fertige Werkstück noch in der Station. Erst nach seiner Entnahme ist
die Station wieder frei. Werkstücke stehen zum Einlegen immer zur Verfügung;
ihre Vorratshaltung wird nicht modelliert.
\begin{enumerate}
  \item Zeichne ein Petrinetz mit drei Stellen für die Zustände der Station und drei Transitionen.
  \item Spiele einen vollständigen Durchlauf an deiner Zeichnung durch.
        Welche Bedeutung hat die Marke, und wo liegt sie nach den einzelnen Schritten?
  \item Warum kann während der Bearbeitung kein zweites Werkstück eingelegt werden?
        Weshalb muss auch ein fertiges Werkstück die Station noch blockieren?
\end{enumerate}

\aufgabetitel{$5$}{Getränkeauswahl}\\
Ein vereinfachter Getränkeautomat wartet anfangs auf eine Münze. Nach dem Einwurf
kann entweder Tee oder Kaffee gewählt oder der Vorgang abgebrochen werden.
Die Wahl gibt sofort das jeweilige Getränk aus; ein Abbruch gibt die Münze zurück.
Danach wartet der Automat wieder auf eine Münze. Solange die Entscheidung aussteht,
kann keine weitere Münze eingeworfen werden. Vorräte und Zubereitungszeiten werden nicht modelliert.
\begin{enumerate}
  \item Modelliere den Ablauf mit zwei Stellen und vier Transitionen.
  \item Welche Transitionen sind unmittelbar nach dem Münzeinwurf aktiviert?
        Begründe, warum mit einer Münze nicht beide Getränke ausgegeben werden können.
  \item Entscheidet das Netz, welche der möglichen Transitionen als Nächstes schaltet?
        Erkläre den Unterschied zwischen ,,aktiviert sein`` und ,,schalten``.
\end{enumerate}

\clearpage
\aufgabetitel{$5$}{Zwei Sensorprüfungen}\\
Beim Start eines Geräts müssen ein Temperatur- und ein Drucksensor geprüft werden.
Anfangs ist das Gerät startbereit. Ein Start stößt beide Prüfungen an.
Die Prüfungen laufen unabhängig voneinander und können in beliebiger Reihenfolge
fertig werden. Erst wenn beide erfolgreich abgeschlossen sind, wird das Gerät
freigegeben. Modelliert wird genau ein Startvorgang ohne Fehlerfälle und ohne Neustart.
\begin{enumerate}
  \item Zeichne ein Petrinetz. Unterscheide für jeden Sensor die Stellen
        ,,Prüfung läuft`` und ,,Prüfung fertig``. Ergänze Stellen für
        ,,startbereit`` und ,,freigegeben``.
  \item Warum können die beiden Prüfungen unabhängig voneinander fertig werden?
        Vergleiche die Verzweigung beim Start mit der Getränkeauswahl aus Aufgabe 2.
  \item Zeige die Marken in deiner Zeichnung, wenn nur die Temperaturprüfung fertig ist.
        Warum kann die Freigabe noch nicht schalten? Was passiert, wenn die Druckprüfung nie fertig wird?
\end{enumerate}

\aufgabetitel{$5$}{Gemeinsamer Arbeitsbereich}\\
Zwei Roboter A und B benutzen denselben Arbeitsbereich. Jeder Roboter befindet
sich entweder außerhalb oder innerhalb dieses Bereichs. Anfangs sind beide
außerhalb und der Bereich ist frei. Ein Roboter darf den Bereich nur betreten,
wenn er frei ist. Beim Verlassen gibt er ihn wieder frei. Beide Roboter können
diesen Ablauf beliebig oft wiederholen.
\begin{enumerate}
  \item Zeichne ein Petrinetz mit je zwei Stellen pro Roboter und einer
        zusätzlichen Stelle ,,Bereich frei``. Verwende je eine Transition
        zum Betreten und zum Verlassen pro Roboter.
  \item Welche Transitionen sind anfangs aktiviert? Welche sind aktiviert,
        nachdem Roboter A den Bereich betreten hat?
  \item Begründe, warum nie beide Roboter gleichzeitig im Bereich sein können.
        Erzwingt das Netz, dass sich A und B abwechseln? Begründe deine Antwort.
\end{enumerate}

\ifloesung{\clearpage}
\begin{loesung}
\textbf{Aufgabe 1: Bearbeitungsstation}

Die Stellen beschreiben Zustände, die Transitionen die Zustandswechsel.
Eine Marke zeigt den aktuellen Zustand der Station an.
\begin{center}
\begin{tikzpicture}[every label/.style=pnlabel]
  \node[pnplace,label=above:{frei}] (f) at (0,0) {}; \pnmark{f}
  \node[pntrans,label=below:{einlegen}] (e) at (2.5,0) {};
  \node[pnplace,label=above:{in Bearbeitung}] (b) at (5,0) {};
  \node[pntrans,label=below:{abschließen}] (a) at (7.5,0) {};
  \node[pnplace,label=above:{fertig}] (g) at (10,0) {};
  \node[pntrans,label=below:{entnehmen}] (n) at (5,-2) {};
  \draw[pnarc] (f)--(e); \draw[pnarc] (e)--(b);
  \draw[pnarc] (b)--(a); \draw[pnarc] (a)--(g);
  \draw[pnarc] (g) |- (n); \draw[pnarc] (n) -| (f);
\end{tikzpicture}
\end{center}

\textbf{Ein Durchlauf:} Anfangs liegt die Marke in ,,frei``. Beim Einlegen
wird sie dort entfernt und in ,,in Bearbeitung`` erzeugt. Beim Abschließen
wechselt sie nach ,,fertig`` und beim Entnehmen zurück nach ,,frei``.
Die Marke steht hier für den Zustand der Station, nicht für einen Werkstückvorrat.

\textbf{Belegung und Reihenfolge:} Während der Bearbeitung fehlt die Marke
in ,,frei``. Deshalb ist ,,einlegen`` nicht aktiviert. Auch ein fertiges Werkstück
belegt die Station noch; erst seine Entnahme ermöglicht das nächste Einlegen.
Würde bereits das Abschließen die Station freigeben, könnte ein neues Werkstück
eingelegt werden, obwohl das alte noch darin liegt.

\textbf{Bedeutung der Marken:} In diesem Netz bleibt genau eine Marke erhalten,
weil jede Transition eine Marke verbraucht und eine erzeugt. Die Station ist
somit immer in genau einem ihrer drei Zustände. Das ist eine Eigenschaft dieses
Modells; Petrinetze können auch mehrere Marken besitzen oder beim Schalten
die Gesamtzahl der Marken verändern (siehe Aufgabe 3).

\clearpage
\textbf{Aufgabe 2: Getränkeauswahl}

Die Stellen unterscheiden ,,warte auf Münze`` und ,,bezahlt``.
Die Getränke-Transitionen fassen jeweils Auswahl und Ausgabe zusammen;
die Abbruchtransition modelliert zugleich die Rückgabe der Münze.
\begin{center}
\begin{tikzpicture}[every label/.style=pnlabel]
  \node[pnplace,label=left:{warte auf Münze}] (w) at (0,0) {}; \pnmark{w}
  \node[pnplace,label=right:{bezahlt}] (z) at (6,0) {};
  \node[pntrans,label=above:{Münze einwerfen}] (e) at (3,1.7) {};
  \node[pntrans,label=above:{Tee wählen/ausgeben}] (t) at (3,0) {};
  \node[pntrans,label=below:{Kaffee wählen/ausgeben}] (k) at (3,-1.5) {};
  \node[pntrans,label=below:{abbrechen/Münze zurück}] (r) at (3,-3) {};
  \draw[pnarc] (w) |- (e); \draw[pnarc] (e) -| (z);
  \draw[pnarc] (z)--(t); \draw[pnarc] (t)--(w);
  \draw[pnarc] (z.south west)--(k.east); \draw[pnarc] (k.west)--(w.south east);
  \draw[pnarc] (z.south) |- (r); \draw[pnarc] (r) -| (w.south);
\end{tikzpicture}
\end{center}

\textbf{Aktivierte Transitionen:} Anfangs trägt nur ,,warte auf Münze`` eine
Marke. Der Münzeinwurf verbraucht diese und erzeugt eine Marke in ,,bezahlt``.
Damit sind Teeauswahl, Kaffeeauswahl und Abbruch aktiviert. Ein weiterer
Münzeinwurf ist in diesem Zustand nicht möglich.

\textbf{Konflikt:} Alle drei Alternativen benötigen dieselbe Marke in ,,bezahlt``.
Schaltet eine davon, verbraucht sie diese Marke. Die anderen beiden sind danach
nicht mehr aktiviert. Eine Münze ermöglicht daher nur eine Getränkeausgabe
oder eine Rückgabe. Anschließend liegt wieder eine Marke in ,,warte auf Münze``.

\textbf{Aktivierung und Auswahl:} ,,Aktiviert`` bedeutet, dass die Voraussetzungen
zum Schalten erfüllt sind. Es bedeutet weder, dass die Transition sofort schaltet,
noch, dass alle aktivierten Transitionen schalten müssen. Das Netz gibt hier
keine Reihenfolge oder Priorität vor. Welche Alternative gewählt wird, hängt
beim realen Automaten von der Bedienung ab; das Modell zeigt die zulässigen Möglichkeiten.

\clearpage
\textbf{Aufgabe 3: Zwei Sensorprüfungen}

Die Starttransition erzeugt je eine Marke für die beiden unabhängigen Prüfungen.
Die Freigabetransition benötigt anschließend beide Fertigmeldungen.
\begin{center}
\begin{tikzpicture}[every label/.style=pnlabel]
  \node[pnplace,label=above:{startbereit}] (s) at (0,0) {}; \pnmark{s}
  \node[pntrans,label=below:{Start}] (a) at (1.7,0) {};
  \node[pnplace,label=above:{Temp. läuft}] (lt) at (3.5,1.6) {};
  \node[pnplace,label=below:{Druck läuft}] (ld) at (3.5,-1.6) {};
  \node[pntrans,label=above:{abschließen}] (t) at (5.6,1.6) {};
  \node[pntrans,label=below:{abschließen}] (d) at (5.6,-1.6) {};
  \node[pnplace,label=above:{Temp. fertig}] (ft) at (7.7,1.6) {};
  \node[pnplace,label=below:{Druck fertig}] (fd) at (7.7,-1.6) {};
  \node[pntrans,label={[yshift=-7mm]below:Freigabe}] (g) at (9.5,0) {};
  \node[pnplace,label=above:{freigegeben}] (b) at (11.3,0) {};
  \draw[pnarc] (s)--(a);
  \draw[pnarc] (a)--(lt); \draw[pnarc] (a)--(ld);
  \draw[pnarc] (lt)--(t); \draw[pnarc] (t)--(ft);
  \draw[pnarc] (ld)--(d); \draw[pnarc] (d)--(fd);
  \draw[pnarc] (ft)--(g); \draw[pnarc] (fd)--(g); \draw[pnarc] (g)--(b);
\end{tikzpicture}
\end{center}
In der Zeichnung steht ,,Temp.`` für Temperatur. Anfangs trägt nur
,,startbereit`` eine Marke; alle anderen Stellen sind leer.

\textbf{Nebenläufigkeit:} Beim Start wird eine Marke verbraucht und je eine Marke
in den beiden Stellen ,,läuft`` erzeugt. Beide Abschluss-Transitionen sind dann
unabhängig aktiviert, da sie Marken aus verschiedenen Stellen benötigen.
Die Temperaturprüfung kann zuerst fertig werden oder die Druckprüfung.
Der Abschluss der einen verhindert den Abschluss der anderen nicht.
Das Netz legt keine Dauer der Prüfungen fest.

\textbf{Unterschied zur Getränkeauswahl:} Dort verzweigen die Pfeile von einer
\emph{Stelle} zu mehreren Transitionen, die um dieselbe Marke konkurrieren.
Hier verzweigen sie von einer \emph{Transition} zu mehreren Stellen:
Beim Start erhalten beide Zweige eine Marke. Es werden also beide Prüfungen
angestoßen, nicht nur eine ausgewählt.

\textbf{Synchronisation:} Ist nur die Temperaturprüfung fertig, liegen Marken
in ,,Temp. fertig`` und ,,Druck läuft``. In ,,Druck fertig`` fehlt noch eine Marke.
Die Freigabe benötigt jedoch eine Marke aus jeder ihrer beiden Eingangsstellen.
Sie kann erst nach Abschluss beider Prüfungen schalten und erzeugt dann eine
Marke in ,,freigegeben``.

Wird die Druckprüfung nie fertig, bleibt die Freigabe aus. Das Modell enthält
weder einen Timeout noch eine Fehlerbehandlung. Nach erfolgreicher Freigabe
kann dagegen absichtlich keine weitere Transition schalten: Der einmalige
Startvorgang ist beendet.

\clearpage
\textbf{Aufgabe 4: Gemeinsamer Arbeitsbereich}

Die Marke in ,,Bereich frei`` ist die Erlaubnis, den gemeinsamen Bereich zu betreten.
Sie wird beim Betreten entnommen und erst beim Verlassen zurückgegeben.
\begin{center}
\begin{tikzpicture}[every label/.style=pnlabel]
  \node[pnplace,label=left:{A außen}] (a0) at (0,2) {}; \pnmark{a0}
  \node[pntrans,label=above:{A betreten}] (ea) at (2.5,2) {};
  \node[pnplace,label=right:{A innen}] (a1) at (5,2) {};
  \node[pntrans,label={[yshift=8mm]above:A verlassen}] (va) at (2.5,3.7) {};
  \node[pnplace,label=above right:{Bereich frei}] (f) at (2.5,0) {}; \pnmark{f}
  \node[pnplace,label=left:{B außen}] (b0) at (0,-2) {}; \pnmark{b0}
  \node[pntrans,label=below:{B betreten}] (eb) at (2.5,-2) {};
  \node[pnplace,label=right:{B innen}] (b1) at (5,-2) {};
  \node[pntrans,label={[yshift=-8mm]below:B verlassen}] (vb) at (2.5,-3.7) {};
  \draw[pnarc] (a0)--(ea); \draw[pnarc] (ea)--(a1);
  \draw[pnarc] (a1) |- (va); \draw[pnarc] (va) -| (a0);
  \draw[pnarc] (b0)--(eb); \draw[pnarc] (eb)--(b1);
  \draw[pnarc] (b1) |- (vb); \draw[pnarc] (vb) -| (b0);
  \draw[pnarc] (f)--(ea); \draw[pnarc] (f)--(eb);
  \draw[pnarc] (va.north)--(2.5,4.5)--(-3,4.5)--(-3,0)--(f.west);
  \draw[pnarc] (vb.south)--(2.5,-4.5)--(8.5,-4.5)--(8.5,0)--(f.east);
\end{tikzpicture}
\end{center}

\textbf{Aktivierte Transitionen:} Anfangs liegt je eine Marke in ,,A außen``,
,,B außen`` und ,,Bereich frei``. Beide Betreten-Transitionen sind aktiviert,
konkurrieren aber um dieselbe Freigabemarke. Nachdem A den Bereich betreten hat,
liegen Marken nur in ,,A innen`` und ,,B außen``. Dann kann ausschließlich
,,A verlassen`` schalten.

\textbf{Gegenseitiger Ausschluss:} Wer den Bereich betritt, verbraucht die
Marke in ,,Bereich frei``. Solange dieser Roboter im Bereich ist, fehlt dem
anderen die erforderliche Marke zum Betreten. Erst das Verlassen gibt die
Erlaubnis zurück. So können nie beide Roboter gleichzeitig im Bereich sein.
Die Marke in ,,Bereich frei`` beschreibt eine verfügbare Ressource; die Marken
in den Roboterstellen beschreiben dagegen deren jeweilige Position.

\textbf{Keine erzwungene Abwechslung:} A darf eintreten, austreten und direkt
erneut eintreten. Nach jedem Verlassen durch A sind wieder beide
Betreten-Transitionen aktiviert; das Netz bevorzugt B nicht. Es garantiert
somit den gegenseitigen Ausschluss, aber keine faire Aufteilung:
B könnte beliebig lange nicht an die Reihe kommen.
\end{loesung}
\end{document}
